\documentclass{article}
\usepackage{amsmath,amssymb,amsthm}
\usepackage{mathpartir}
\usepackage{stmaryrd}
\usepackage{xcolor}

\newtheorem{theorem}{Theorem}[section]
\newtheorem{lemma}[theorem]{Lemma}
\newtheorem{definition}[theorem]{Definition}
\newtheorem{corollary}[theorem]{Corollary}

\title{Sounio: A Type System with Semantic Distance\\
\large Formal Sketch (not mechanized)}
\author{Demetrios Chiuratto Agourakis}
\date{December 2024}

\begin{document}

\maketitle

\begin{abstract}
This document is a formal sketch, not a mechanized proof or current compiler
contract. It presents an intended type system where types may be ontological
terms and compatibility may be determined by semantic distance. The syntax,
typing rules, and proof sketches below should be cited as design material
unless an exact Lean artifact or executable gate is linked for the specific
claim.
\end{abstract}

\section*{Status}

\textbf{FORMAL SKETCH --- not mechanized.} The theorem and proof environments
below are retained as mathematical design notes. They are not, by themselves,
evidence that the current compiler implements semantic-distance subtyping or
that the metric properties are mechanized.

\section{Semantic Metric Types}

\subsection{Syntax}

\begin{definition}[Type Syntax]
\begin{align*}
\text{Types } \tau &::= i \mid \tau_1 \to \tau_2 \mid \forall \alpha. \tau \mid \text{Knowledge}[\tau, \epsilon, \delta, \Phi] \\
\text{IRIs } i &::= \text{CURIE} \mid \text{URI} \\
\text{Distance } d &\in [0, 1] \\
\text{Confidence } c &\in [0, 1] \\
\text{Threshold } \theta &\in [0, 1] \\
\text{Provenance } \Phi &::= \text{Source} \mid \Phi_1 \cdot \Phi_2 \mid \text{Coerce}(d)
\end{align*}
\end{definition}

\begin{definition}[Expression Syntax]
\begin{align*}
\text{Expressions } e &::= x \mid \lambda x:\tau. e \mid e_1 \, e_2 \mid e : \tau \\
                      &\mid \text{let } x = e_1 \text{ in } e_2 \\
                      &\mid \text{Knowledge}(e, c, \delta, \Phi) \\
                      &\mid e \text{ as } \tau
\end{align*}
\end{definition}

\subsection{Semantic Distance Function}

\begin{definition}[Semantic Distance]
We define $d : \text{IRI} \times \text{IRI} \to [0, 1]$ as:
\[
d(i, j) = w_p \cdot d_{\text{path}}(i, j) + w_c \cdot d_{\text{IC}}(i, j) + w_e \cdot d_{\text{emb}}(i, j)
\]
where $w_p + w_c + w_e = 1$ and:
\begin{itemize}
\item $d_{\text{path}}(i, j) = \frac{\text{path\_length}(i, j)}{2 \cdot \text{max\_depth}}$ is normalized path distance
\item $d_{\text{IC}}(i, j) = 1 - \frac{2 \cdot \text{IC}(\text{LCA}(i,j))}{\text{IC}(i) + \text{IC}(j)}$ is information content distance
\item $d_{\text{emb}}(i, j) = \frac{1 - \cos(\vec{v}_i, \vec{v}_j)}{2}$ is embedding distance
\end{itemize}
\end{definition}

\begin{theorem}[Metric Properties (sketch)]
The distance function $d$ satisfies:
\begin{enumerate}
\item $d(i, i) = 0$ \hfill (Identity)
\item $d(i, j) = d(j, i)$ \hfill (Symmetry)
\item $d(i, k) \leq d(i, j) + d(j, k)$ \hfill (Triangle inequality)
\end{enumerate}
\end{theorem}

\begin{proof}
Each component satisfies metric properties:
\begin{enumerate}
\item Path distance: shortest path from $i$ to $i$ is 0.
\item IC distance: When $i = j$, LCA$(i,i) = i$, so numerator equals denominator.
\item Embedding distance: $\cos(\vec{v}_i, \vec{v}_i) = 1$, so distance is 0.
\end{enumerate}
Symmetry follows from symmetry of each component. Triangle inequality follows
from the triangle inequality of each component and convex combination.
\end{proof}

\subsection{Subtyping with Distance}

Traditional subtyping: $\tau_1 <: \tau_2$ (boolean)

Semantic subtyping: $\tau_1 <:_d \tau_2$ where $d = d(\tau_1, \tau_2)$

\begin{mathpar}
\inferrule[S-Refl]
  { }
  {\tau <:_0 \tau}

\inferrule[S-Onto-Sub]
  {i_1 \sqsubseteq_{\mathcal{O}} i_2 \\ d = d_{\text{path}}(i_1, i_2)}
  {i_1 <:_d i_2}

\inferrule[S-Onto-Prox]
  {d(i_1, i_2) = d \\ d \leq \theta_{\text{max}}}
  {i_1 <:_d i_2}

\inferrule[S-Trans]
  {\tau_1 <:_{d_1} \tau_2 \\ \tau_2 <:_{d_2} \tau_3}
  {\tau_1 <:_{\min(d_1 + d_2, 1)} \tau_3}

\inferrule[S-Arrow]
  {\tau_1' <:_{d_1} \tau_1 \\ \tau_2 <:_{d_2} \tau_2'}
  {\tau_1 \to \tau_2 <:_{\max(d_1, d_2)} \tau_1' \to \tau_2'}

\inferrule[S-Knowledge]
  {\tau <:_d \tau' \\ c' = c \cdot (1 - \alpha \cdot d)}
  {\text{Knowledge}[\tau, c, \delta, \Phi] <:_d \text{Knowledge}[\tau', c', \delta, \Phi \cdot \text{Coerce}(d)]}
\end{mathpar}

\subsection{Typing Rules}

\begin{definition}[Typing Judgment]
We write $\Gamma \vdash e : \tau \triangleright d$ to mean that under context
$\Gamma$, expression $e$ has type $\tau$ with accumulated semantic distance $d$.
\end{definition}

\begin{mathpar}
\inferrule[T-Var]
  {x : \tau \in \Gamma}
  {\Gamma \vdash x : \tau \triangleright 0}

\inferrule[T-Abs]
  {\Gamma, x : \tau_1 \vdash e : \tau_2 \triangleright d}
  {\Gamma \vdash \lambda x:\tau_1. e : \tau_1 \to \tau_2 \triangleright d}

\inferrule[T-App]
  {\Gamma \vdash e_1 : \tau_1 \to \tau_2 \triangleright d_1 \\
   \Gamma \vdash e_2 : \tau_1' \triangleright d_2 \\
   \tau_1' <:_{d_s} \tau_1 \\
   d_s \leq \theta_{\text{implicit}}}
  {\Gamma \vdash e_1 \, e_2 : \tau_2 \triangleright d_1 + d_2 + d_s}

\inferrule[T-App-Explicit]
  {\Gamma \vdash e_1 : \tau_1 \to \tau_2 \triangleright d_1 \\
   \Gamma \vdash e_2 : \tau_1' \triangleright d_2 \\
   \tau_1' <:_{d_s} \tau_1 \\
   \theta_{\text{implicit}} < d_s \leq \theta_{\text{explicit}}}
  {\Gamma \vdash e_1 \, (e_2 \text{ as } \tau_1) : \tau_2 \triangleright d_1 + d_2 + d_s}

\inferrule[T-Let]
  {\Gamma \vdash e_1 : \tau_1 \triangleright d_1 \\
   \Gamma, x : \tau_1 \vdash e_2 : \tau_2 \triangleright d_2}
  {\Gamma \vdash \text{let } x = e_1 \text{ in } e_2 : \tau_2 \triangleright d_1 + d_2}

\inferrule[T-Knowledge]
  {\Gamma \vdash e : \tau \triangleright d \\
   c \in [0, 1] \\
   \delta : \text{Domain} \\
   \Phi : \text{Provenance}}
  {\Gamma \vdash \text{Knowledge}(e, c, \delta, \Phi) : \text{Knowledge}[\tau, c, \delta, \Phi] \triangleright d}

\inferrule[T-Coerce]
  {\Gamma \vdash e : \text{Knowledge}[\tau_1, c, \delta, \Phi] \triangleright d \\
   \tau_1 <:_{d_s} \tau_2 \\
   c' = c \cdot (1 - \alpha \cdot d_s)}
  {\Gamma \vdash e \text{ as } \tau_2 : \text{Knowledge}[\tau_2, c', \delta, \Phi \cdot \text{Coerce}(d_s)] \triangleright d + d_s}

\inferrule[T-Sub]
  {\Gamma \vdash e : \tau \triangleright d \\
   \tau <:_{d_s} \tau' \\
   d_s \leq \theta_{\text{implicit}}}
  {\Gamma \vdash e : \tau' \triangleright d + d_s}
\end{mathpar}

\subsection{Confidence Propagation}

\begin{definition}[Confidence Degradation]
When types are coerced across semantic distance $d$, confidence degrades:
\[
c' = c \cdot (1 - \alpha \cdot d)
\]
where $\alpha \in [0, 1]$ is the degradation factor (typically 0.1--0.2).
\end{definition}

\begin{lemma}[Confidence Bounds]
For any sequence of coercions with distances $d_1, \ldots, d_n$:
\[
c_n = c_0 \cdot \prod_{i=1}^{n} (1 - \alpha \cdot d_i) \geq c_0 \cdot (1 - \alpha)^n
\]
\end{lemma}

\section{Knowledge Type System}

The $\text{Knowledge}[\tau, \epsilon, \delta, \Phi]$ type encapsulates epistemic values with:
\begin{itemize}
\item $\tau$: The underlying value type
\item $\epsilon$: Uncertainty (standard deviation or interval half-width)
\item $\delta$: Domain/ontology context
\item $\Phi$: Provenance chain tracking measurement sources and transformations
\end{itemize}

\subsection{Knowledge Type Syntax}

\begin{definition}[Extended Expression Syntax for Knowledge]
\begin{align*}
e &::= \ldots \\
  &\mid \text{knowledge}(e_v, e_\epsilon, \delta, \Phi) \tag{Introduction} \\
  &\mid \text{unwrap}(e) \tag{Elimination} \\
  &\mid \text{unwrap\_checked}(e, c_{\min}) \tag{Checked Elimination} \\
  &\mid e_1 \oplus_K e_2 \tag{Binary Operations} \\
  &\mid \text{combine}(\overline{e}, w) \tag{Fusion}
\end{align*}
where $\oplus \in \{+, -, \times, \div\}$ and $w$ is a weight vector.
\end{definition}

\subsection{Knowledge Typing Rules}

\begin{mathpar}
\inferrule[T-Knowledge-Intro]
  {\Gamma \vdash e_v : \tau \triangleright d_v \\
   \Gamma \vdash e_\epsilon : \mathbb{R}_{\geq 0} \triangleright d_\epsilon \\
   \delta : \text{Domain} \\
   \Phi : \text{Provenance} \\
   c = \text{conf}(\Phi)}
  {\Gamma \vdash \text{knowledge}(e_v, e_\epsilon, \delta, \Phi) :
   \text{Knowledge}[\tau, \epsilon, \delta, \Phi] \triangleright d_v + d_\epsilon}
\end{mathpar}

The introduction rule constructs a Knowledge value from a base value $e_v$,
an uncertainty $e_\epsilon$, a domain context $\delta$, and provenance $\Phi$.
The confidence $c$ is derived from the provenance chain.

\begin{mathpar}
\inferrule[T-Knowledge-Elim]
  {\Gamma \vdash e : \text{Knowledge}[\tau, \epsilon, \delta, \Phi] \triangleright d \\
   c = \text{conf}(\Phi)}
  {\Gamma \vdash \text{unwrap}(e) : \tau \triangleright d}

\inferrule[T-Knowledge-Elim-Checked]
  {\Gamma \vdash e : \text{Knowledge}[\tau, \epsilon, \delta, \Phi] \triangleright d \\
   c = \text{conf}(\Phi) \\
   c \geq c_{\min}}
  {\Gamma \vdash \text{unwrap\_checked}(e, c_{\min}) : \tau \triangleright d}
\end{mathpar}

The basic elimination rule extracts the underlying value, discarding epistemic metadata.
The checked variant enforces a minimum confidence threshold $c_{\min}$, failing at
compile time if the static confidence is insufficient or at runtime if dynamic
confidence drops below the threshold.

\begin{mathpar}
\inferrule[T-Knowledge-Coerce]
  {\Gamma \vdash e : \text{Knowledge}[\tau_1, \epsilon, \delta_1, \Phi] \triangleright d \\\\
   \tau_1 <:_{d_s} \tau_2 \\
   d_s \leq \theta_{\text{explicit}} \\\\
   c = \text{conf}(\Phi) \\
   c' = c \cdot (1 - \alpha \cdot d_s) \\\\
   \Phi' = \Phi \cdot \text{Coerce}(\tau_1, \tau_2, d_s)}
  {\Gamma \vdash e \text{ as } \text{Knowledge}[\tau_2] :
   \text{Knowledge}[\tau_2, \epsilon, \delta_2, \Phi'] \triangleright d + d_s}
\end{mathpar}

Cross-ontology coercion applies the confidence degradation formula:
\[
c' = c \cdot (1 - \alpha \cdot d_s)
\]
where $\alpha$ is the degradation factor and $d_s$ is the semantic distance
between source and target types. The provenance chain records the coercion
for auditability.

\subsection{GUM Uncertainty Propagation}

Following the ISO Guide to the Expression of Uncertainty in Measurement (GUM),
we propagate uncertainties through arithmetic operations assuming uncorrelated inputs.

\begin{definition}[GUM Variance Propagation]
For a function $f(x_1, \ldots, x_n)$ with inputs having uncertainties $\epsilon_1, \ldots, \epsilon_n$:
\[
\epsilon_f^2 = \sum_{i=1}^{n} \left( \frac{\partial f}{\partial x_i} \right)^2 \epsilon_i^2
\]
\end{definition}

For binary arithmetic operations on Knowledge values:

\begin{definition}[Binary Operation Uncertainty Rules]
Let $K_1 = \text{Knowledge}[\tau, \epsilon_1, \delta, \Phi_1]$ with value $v_1$ and
$K_2 = \text{Knowledge}[\tau, \epsilon_2, \delta, \Phi_2]$ with value $v_2$:

\begin{align*}
\text{Addition/Subtraction:} \quad & \epsilon_{v_1 \pm v_2}^2 = \epsilon_1^2 + \epsilon_2^2 \\[0.5em]
\text{Multiplication:} \quad & \epsilon_{v_1 \cdot v_2}^2 = v_2^2 \epsilon_1^2 + v_1^2 \epsilon_2^2 \\[0.5em]
\text{Division:} \quad & \epsilon_{v_1 / v_2}^2 = \frac{1}{v_2^2} \epsilon_1^2 + \frac{v_1^2}{v_2^4} \epsilon_2^2
\end{align*}

In relative uncertainty form (coefficient of variation):
\begin{align*}
\text{Multiplication:} \quad & \left(\frac{\epsilon_{v_1 \cdot v_2}}{v_1 \cdot v_2}\right)^2 = \left(\frac{\epsilon_1}{v_1}\right)^2 + \left(\frac{\epsilon_2}{v_2}\right)^2 \\[0.5em]
\text{Division:} \quad & \left(\frac{\epsilon_{v_1 / v_2}}{v_1 / v_2}\right)^2 = \left(\frac{\epsilon_1}{v_1}\right)^2 + \left(\frac{\epsilon_2}{v_2}\right)^2
\end{align*}
\end{definition}

\begin{mathpar}
\inferrule[T-Knowledge-Binary-Add]
  {\Gamma \vdash e_1 : \text{Knowledge}[\tau, \epsilon_1, \delta, \Phi_1] \triangleright d_1 \\\\
   \Gamma \vdash e_2 : \text{Knowledge}[\tau, \epsilon_2, \delta, \Phi_2] \triangleright d_2 \\\\
   \epsilon' = \sqrt{\epsilon_1^2 + \epsilon_2^2} \\
   c' = \min(\text{conf}(\Phi_1), \text{conf}(\Phi_2)) \\\\
   \Phi' = \text{BinOp}(+, \Phi_1, \Phi_2)}
  {\Gamma \vdash e_1 +_K e_2 : \text{Knowledge}[\tau, \epsilon', \delta, \Phi'] \triangleright d_1 + d_2}

\inferrule[T-Knowledge-Binary-Sub]
  {\Gamma \vdash e_1 : \text{Knowledge}[\tau, \epsilon_1, \delta, \Phi_1] \triangleright d_1 \\\\
   \Gamma \vdash e_2 : \text{Knowledge}[\tau, \epsilon_2, \delta, \Phi_2] \triangleright d_2 \\\\
   \epsilon' = \sqrt{\epsilon_1^2 + \epsilon_2^2} \\
   c' = \min(\text{conf}(\Phi_1), \text{conf}(\Phi_2)) \\\\
   \Phi' = \text{BinOp}(-, \Phi_1, \Phi_2)}
  {\Gamma \vdash e_1 -_K e_2 : \text{Knowledge}[\tau, \epsilon', \delta, \Phi'] \triangleright d_1 + d_2}

\inferrule[T-Knowledge-Binary-Mul]
  {\Gamma \vdash e_1 : \text{Knowledge}[\tau_1, \epsilon_1, \delta, \Phi_1] \triangleright d_1 \\
   v_1 = \text{val}(e_1) \\\\
   \Gamma \vdash e_2 : \text{Knowledge}[\tau_2, \epsilon_2, \delta, \Phi_2] \triangleright d_2 \\
   v_2 = \text{val}(e_2) \\\\
   \tau' = \tau_1 \cdot \tau_2 \\
   \epsilon' = \sqrt{v_2^2 \epsilon_1^2 + v_1^2 \epsilon_2^2} \\\\
   c' = \min(\text{conf}(\Phi_1), \text{conf}(\Phi_2)) \\
   \Phi' = \text{BinOp}(\times, \Phi_1, \Phi_2)}
  {\Gamma \vdash e_1 \times_K e_2 : \text{Knowledge}[\tau', \epsilon', \delta, \Phi'] \triangleright d_1 + d_2}

\inferrule[T-Knowledge-Binary-Div]
  {\Gamma \vdash e_1 : \text{Knowledge}[\tau_1, \epsilon_1, \delta, \Phi_1] \triangleright d_1 \\
   v_1 = \text{val}(e_1) \\\\
   \Gamma \vdash e_2 : \text{Knowledge}[\tau_2, \epsilon_2, \delta, \Phi_2] \triangleright d_2 \\
   v_2 = \text{val}(e_2) \\
   v_2 \neq 0 \\\\
   \tau' = \tau_1 / \tau_2 \\
   \epsilon' = \sqrt{\frac{\epsilon_1^2}{v_2^2} + \frac{v_1^2 \epsilon_2^2}{v_2^4}} \\\\
   c' = \min(\text{conf}(\Phi_1), \text{conf}(\Phi_2)) \\
   \Phi' = \text{BinOp}(\div, \Phi_1, \Phi_2)}
  {\Gamma \vdash e_1 \div_K e_2 : \text{Knowledge}[\tau', \epsilon', \delta, \Phi'] \triangleright d_1 + d_2}
\end{mathpar}

\subsection{Knowledge Combination}

Multiple Knowledge sources can be fused using weighted combination,
typically for sensor fusion or meta-analysis.

\begin{definition}[Weighted Combination]
Given $n$ Knowledge values $K_i = \text{Knowledge}[\tau, \epsilon_i, \delta, \Phi_i]$
with values $v_i$ and weights $w_i$ (where $\sum_i w_i = 1$):

\begin{align*}
v_{\text{combined}} &= \sum_{i=1}^{n} w_i \cdot v_i \\[0.5em]
\epsilon_{\text{combined}}^2 &= \sum_{i=1}^{n} w_i^2 \cdot \epsilon_i^2 \\[0.5em]
c_{\text{combined}} &= \min_{i} \text{conf}(\Phi_i) \cdot \left(1 - \frac{\text{Var}(\{v_i\})}{\text{scale}}\right)
\end{align*}

The confidence term penalizes disagreement among sources via their variance.
\end{definition}

For optimal combination under the assumption of independent Gaussian errors,
the minimum-variance unbiased estimator uses inverse-variance weighting:

\begin{definition}[Optimal Weights]
\[
w_i^* = \frac{1/\epsilon_i^2}{\sum_j 1/\epsilon_j^2}
\]
yielding combined uncertainty:
\[
\epsilon_{\text{combined}}^2 = \frac{1}{\sum_i 1/\epsilon_i^2}
\]
\end{definition}

\begin{mathpar}
\inferrule[T-Knowledge-Combine]
  {\forall i \in 1..n. \; \Gamma \vdash e_i : \text{Knowledge}[\tau, \epsilon_i, \delta, \Phi_i] \triangleright d_i \\\\
   w = (w_1, \ldots, w_n) \\
   \sum_i w_i = 1 \\\\
   \epsilon' = \sqrt{\sum_i w_i^2 \cdot \epsilon_i^2} \\
   c' = \text{combine\_conf}(\overline{\Phi}, \overline{v}, w) \\\\
   \Phi' = \text{Combine}(\overline{\Phi}, w)}
  {\Gamma \vdash \text{combine}(\overline{e}, w) :
   \text{Knowledge}[\tau, \epsilon', \delta, \Phi'] \triangleright \sum_i d_i}
\end{mathpar}

\subsection{Correlated Uncertainty Propagation}

When Knowledge values share common sources (detectable via provenance intersection),
the covariance must be included:

\begin{definition}[Correlated GUM Propagation]
For correlated inputs with covariance matrix $\Sigma$:
\[
\epsilon_f^2 = \mathbf{J}^T \Sigma \mathbf{J}
\]
where $\mathbf{J}$ is the Jacobian vector $(\partial f / \partial x_i)$.

For addition with correlation coefficient $\rho$:
\[
\epsilon_{v_1 + v_2}^2 = \epsilon_1^2 + \epsilon_2^2 + 2\rho\epsilon_1\epsilon_2
\]
\end{definition}

\begin{mathpar}
\inferrule[T-Knowledge-Binary-Correlated]
  {\Gamma \vdash e_1 : \text{Knowledge}[\tau, \epsilon_1, \delta, \Phi_1] \triangleright d_1 \\\\
   \Gamma \vdash e_2 : \text{Knowledge}[\tau, \epsilon_2, \delta, \Phi_2] \triangleright d_2 \\\\
   \rho = \text{correlation}(\Phi_1, \Phi_2) \\
   \epsilon' = \sqrt{\epsilon_1^2 + \epsilon_2^2 + 2\rho\epsilon_1\epsilon_2} \\\\
   \Phi' = \text{BinOp}(+, \Phi_1, \Phi_2, \rho)}
  {\Gamma \vdash e_1 +_K^{\text{corr}} e_2 : \text{Knowledge}[\tau, \epsilon', \delta, \Phi'] \triangleright d_1 + d_2}
\end{mathpar}

\section{Soundness}

\begin{theorem}[Type Safety]
If $\Gamma \vdash e : \tau \triangleright d$ and $e \to^* v$, then
$\Gamma \vdash v : \tau' \triangleright d'$ where $\tau' <:_{d''} \tau$
and $d' \leq d$.
\end{theorem}

\begin{proof}[Proof Sketch]
By induction on the derivation of $\Gamma \vdash e : \tau \triangleright d$
and the evaluation steps. The key insight is that each evaluation step
either preserves the type exactly or reduces to a subtype with accumulated
distance. The distance tracking ensures we never exceed the original
type's compatibility bound.
\end{proof}

\begin{theorem}[Distance Monotonicity]
If $\Gamma \vdash e_1 : \tau \triangleright d_1$ and $e_1 \to e_2$,
then $\Gamma \vdash e_2 : \tau \triangleright d_2$ where $d_2 \leq d_1$.
\end{theorem}

\begin{proof}[Proof Sketch]
Evaluation does not introduce new type coercions. Each step either:
\begin{enumerate}
\item Substitutes values of compatible types (distance preserved)
\item Reduces applications (distance from source preserved)
\item Eliminates let bindings (distance combined but not increased)
\end{enumerate}
\end{proof}

\begin{theorem}[Confidence Preservation]
For any well-typed computation $e \to^* v$ where $e : \text{Knowledge}[\tau, c, \delta, \Phi]$:
\[
\text{conf}(v) \leq c
\]
Confidence is non-increasing through computation.
\end{theorem}

\begin{proof}
Confidence only changes through explicit coercion operations (T-Coerce),
each of which multiplies by $(1 - \alpha \cdot d) \leq 1$. Pure computation
preserves confidence exactly.
\end{proof}

\begin{corollary}[Epistemic Soundness]
If a value $v$ has $\text{Knowledge}[\tau, c, \delta, \Phi]$ type at runtime,
then the confidence $c$ accurately reflects the accumulated uncertainty from
all type coercions in its derivation, as recorded in provenance $\Phi$.
\end{corollary}

\subsection{Knowledge Type Soundness}

\begin{theorem}[Knowledge Type Soundness]
\label{thm:knowledge-soundness}
The Knowledge type system satisfies the following properties:

\begin{enumerate}
\item \textbf{(Uncertainty Soundness)} For any well-typed arithmetic expression
$e : \text{Knowledge}[\tau, \epsilon, \delta, \Phi]$ constructed from inputs
$K_1, \ldots, K_n$ with true uncertainties $\sigma_1, \ldots, \sigma_n$:
\[
\epsilon \geq \epsilon_{\text{GUM}}
\]
where $\epsilon_{\text{GUM}}$ is the uncertainty computed by GUM propagation.
The computed uncertainty is a conservative (over-)estimate.

\item \textbf{(Confidence Monotonicity)} For any derivation $D$ of
$\Gamma \vdash e : \text{Knowledge}[\tau, \epsilon, \delta, \Phi]$:
\[
\text{conf}(\Phi) \leq \min_{K_i \in \text{sources}(e)} \text{conf}(\Phi_i)
\]
Confidence cannot increase through computation.

\item \textbf{(Coercion Tracking)} Every coercion in a derivation is recorded
in the provenance chain $\Phi$:
\[
\#\text{Coerce}(\Phi) = \#\{\text{applications of T-Knowledge-Coerce in } D\}
\]

\item \textbf{(Degradation Bound)} For a Knowledge value with $n$ coercions
at distances $d_1, \ldots, d_n$:
\[
\text{conf}(\Phi) = c_0 \cdot \prod_{i=1}^{n} (1 - \alpha \cdot d_i)
\]
where $c_0$ is the initial source confidence.
\end{enumerate}
\end{theorem}

\begin{proof}
\textbf{(1)} By structural induction on the derivation. Base case: T-Knowledge-Intro
assigns the declared uncertainty. Inductive case: each binary operation rule
(T-Knowledge-Binary-*) computes uncertainty via GUM formulas, which are exact
for linear operations and first-order approximations for nonlinear operations.
The approximation is conservative when higher-order terms are positive.

\textbf{(2)} By induction. T-Knowledge-Intro establishes initial confidence from
provenance. T-Knowledge-Binary-* rules take $\min$ of input confidences.
T-Knowledge-Coerce multiplies by $(1 - \alpha \cdot d) \leq 1$.
T-Knowledge-Combine uses a penalized minimum. All operations are non-increasing.

\textbf{(3)} By construction: T-Knowledge-Coerce appends $\text{Coerce}(\cdot)$
to $\Phi$, and no other rule adds coercion records. Provenance concatenation
is associative and injective.

\textbf{(4)} Direct consequence of the T-Knowledge-Coerce rule applied $n$ times,
with confidence degradation formula $c' = c \cdot (1 - \alpha \cdot d)$
composed multiplicatively.
\end{proof}

\begin{theorem}[GUM Compliance]
\label{thm:gum-compliance}
The uncertainty propagation rules in the Knowledge type system are compliant
with ISO/IEC Guide 98-3:2008 (GUM) for:
\begin{enumerate}
\item Uncorrelated Type A uncertainties (standard deviations)
\item Linear measurement models
\item First-order Taylor expansion for nonlinear models
\end{enumerate}
\end{theorem}

\begin{proof}
The binary operation rules directly implement GUM equations:
\begin{itemize}
\item GUM Equation 10: $u_c^2(y) = \sum_i (c_i)^2 u^2(x_i)$ where $c_i = \partial f / \partial x_i$
\item For $f = x_1 + x_2$: $c_1 = c_2 = 1$, giving $\epsilon^2 = \epsilon_1^2 + \epsilon_2^2$
\item For $f = x_1 \cdot x_2$: $c_1 = x_2$, $c_2 = x_1$, giving $\epsilon^2 = x_2^2\epsilon_1^2 + x_1^2\epsilon_2^2$
\item For $f = x_1 / x_2$: $c_1 = 1/x_2$, $c_2 = -x_1/x_2^2$, giving the stated formula
\end{itemize}
The T-Knowledge-Binary-Correlated rule implements GUM Equation 13 for correlated inputs.
\end{proof}

\section{Disjointness}

\begin{definition}[Disjoint Types]
Types $\tau_1$ and $\tau_2$ are disjoint, written $\tau_1 \perp \tau_2$,
if the ontology declares them as disjoint classes.
\end{definition}

\begin{mathpar}
\inferrule[Disjoint-Error]
  {\tau_1 \perp \tau_2}
  {\tau_1 \not<:_d \tau_2 \text{ for any } d}
\end{mathpar}

\begin{theorem}[Disjoint Safety]
If $\tau_1 \perp \tau_2$ in the ontology, then no well-typed program can
coerce a value of type $\tau_1$ to type $\tau_2$.
\end{theorem}

\section{Cross-Ontology Compatibility}

\begin{definition}[SSSOM Mapping]
A mapping $m : i_1 \leftrightarrow i_2$ with predicate $p$ and confidence $c_m$
relates IRIs across ontologies.
\end{definition}

\begin{mathpar}
\inferrule[S-Cross-Onto]
  {m : i_1 \leftrightarrow i_2 \\
   p \in \{\text{exactMatch}, \text{closeMatch}\} \\
   c_m \geq \theta_{\text{mapping}}}
  {i_1 <:_{1 - c_m} i_2}
\end{mathpar}

\section{Implementation Notes}

\subsection{Threshold Configuration}

\begin{itemize}
\item $\theta_{\text{implicit}} = 0.15$: Maximum distance for implicit coercion
\item $\theta_{\text{explicit}} = 0.40$: Maximum distance for explicit coercion
\item $\theta_{\text{max}} = 0.60$: Maximum distance for any coercion
\item $\alpha = 0.15$: Confidence degradation factor
\end{itemize}

\subsection{Distance Weight Configuration}

Default weights for multi-modal distance:
\begin{itemize}
\item $w_p = 0.35$: Path distance weight
\item $w_c = 0.25$: Information content weight
\item $w_e = 0.40$: Embedding distance weight
\end{itemize}

\end{document}
